%% hopfion-framework/gravity/dark_sector.tex
%% Sources: main_paper7.tex
%% Cosmological constant, dark energy w_a, dark matter c_s, BAO, Vainshtein
%% Auto-generated from project files. Edit source papers to update.


%════════════════════════════════════════════════════════════════════
% Cosmological Constant and Dark Energy — Paper VII
%════════════════════════════════════════════════════════════════════
The field equation of~\eqref{P7:eq:parent_action} at homogeneous large
$\rho$ has the self-consistent attractor solution $\rho=\rinfty$ with
\begin{equation}
  \rinfty = \frac{\vp}{\beta},
  \label{P7:eq:rho_inf}
\end{equation}
giving $\rinfty\approx16.18\,\MPl^3$ for $\beta=0.1$. The cosmological
constant today is
\begin{equation}
  \boxed{\Lobs = \Lambda_0\,e^{-\vp^6\rinfty}
               = \Lambda_0\,e^{-\vp^7/\beta}
               \approx \MPl^4\times10^{-122}.}
  \label{P7:eq:Lambda_obs}
\end{equation}
The exponent's normalisation constant $c=\rho/\kappa$ between the
$k$-essence field and the Faddeev--Niemi gradient density is
\begin{equation}
  c \;=\; \vp^3 + p_0^2 \;=\; \vp^3\Bigl(1+\frac{1}{\vp^7}\Bigr),
  \label{P7:eq:cc_c}
\end{equation}
where $p_0=1/(1+\bst\rinfty)=1/\vp^2$ is the vacuum Born weight, giving
$\approx10^{-122}\,\MPl^4$ in agreement with observation.

\begin{proposition}[Cosmological-constant normalisation]
\label{P7:prop:cc_norm}
On the density-feedback Euler--Lagrange manifold, the normalisation
between the $k$-essence field $\rho$ and the Faddeev--Niemi gradient
density $\kappa$ is
\begin{equation}
  c \;\equiv\; \frac{\rho}{\kappa} \;=\; \vp^3 + p_0^{\,Q_H},
  \label{P7:eq:cc_norm}
\end{equation}
with $p_0 = 1/V^{*2}$ the WZW vacuum weight, $V^* = S_{0,1/2}/S_{0,0}$ the
virial value, and $Q_H$ the Hopf charge. For the $Q_H=2$ icosahedral
condensate $c = \vp^3 + \vp^{-4}$, giving
$\Lobs = \Lambda_0 e^{-\vp^7 c/\bst}\approx10^{-122}\,\MPl^4$.
\end{proposition}
\status{published}
\source{P7:prop:cc_norm}
\begin{proof}
\emph{Tree normalisation.} The parent kinetic density
$|\nabla\rho|^2/[\vp^6(1+\bst\kappa)]$, $\kappa=|\nabla\rho|^2$,
canonically normalises via $d\chi/d\rho=[\vp^6(1+\bst\kappa)]^{-1/2}$;
in the ultraviolet limit $\bst\kappa\to0$ this is $\vp^{-3}$, so
$c_{\mathrm{tree}}=\vp^3=\sqrt\lam$.

\emph{Vacuum weight.} The $k=3$ modular $S$-matrix satisfies
$S_{1/2,1/2}=S_{0,0}$, so the vacuum monodromy scalar is
$p_0=S_{1/2,1/2}S_{0,0}/S_{0,1/2}^2=(S_{0,0}/S_{0,1/2})^2=1/V^{*2}$,
with $V^*=S_{0,1/2}/S_{0,0}=\vp$ proved for all $R_0$ by three-fermion
Witten bosonization (Paper~III). Hence $p_0=1/\vp^2$.

\emph{Linking count.} The vacuum-dressing correction to $c$ is the
field self-energy: the Hopf charge is the linking number of the
field's fibre preimages, so the self-energy loop links the condensate
flux $Q_H$ times and $c_{\mathrm{vac}}=p_0^{\,Q_H}$.

\emph{Assembly.} $c=c_{\mathrm{tree}}+c_{\mathrm{vac}}=\vp^3+p_0^{\,Q_H}$.
At $Q_H=2$, $c=\vp^3+\vp^{-4}=4.382$, whence $\beta=\bst/c=0.1032$,
$\rinfty=\vp/\beta=15.68$, and $\Lobs=\Lambda_0e^{-\vp^6\rinfty}=10^{-122.2}\,\MPl^4$.
\end{proof}
\status{published}
\source{P7:prop:cc_norm}

\begin{remark}[Status of the linking step]
\label{P7:rem:cc_norm_status}
The tree normalisation, the vacuum weight $p_0=1/V^{*2}$ (with
$V^*=\vp$ proved by bosonization), and the self-energy Chern--Simons
Wilson-loop evaluation are all computed. The one topological input is
that the self-energy loop's linking with the condensate flux equals
the Hopf charge --- the definition of the Hopf invariant. This is the
\emph{geometric} Hopf charge $Q_H=2$, distinct from the
\emph{algebraic} charge $Q=10$ that organises the geometric
quantisation (Paper~X).
\end{remark}
\status{published}
\source{P7:rem:cc_norm_status}

When $\rho$ is frozen at $\rinfty$ ($\dot\rho=0$), the equation of
state is $w=-1$ exactly, reproducing $\Lambda$CDM.
The dark energy scalar mass at the potential minimum is
\begin{equation}
  m_\xi^2 = \vp^{20}\,\Lobs\,\MPl^4,
  \label{P7:eq:mxi}
\end{equation}
in the Jordan frame. The physical mass governing the cosmological
(thawing) dynamics is the Einstein-frame value, obtained after the
conformal rescaling by $F(\rinfty)$~\eqref{P7:eq:Feff_derived}:
because that nonminimal coupling is strong ($1+3\vp^7\approx88$), the
conformal kinetic normalisation cancels the $\vp^{20}$ curvature,
leaving the $\vp$-power-free Hubble-scale result
$m_\xi\sim\sqrt{\Lobs}/\MPl\approx H_0\approx10^{-33}\,\mathrm{eV}$.
The Vainshtein/screening analysis uses the Jordan-frame
expression~\eqref{P7:eq:mxi} directly; the resulting PPN parameter
resolves Cassini, and the outcome is unchanged if the physical mass is
used, since $r_V$ remains far larger than the solar system either way.

\begin{remark}[Dark-energy sound speed]
\label{P7:rem:cs_de}
The sound speed of the dark-energy scalar's perturbations at the
attractor is fixed with no free parameters. With
$K(\rho)=1/[\vp^6(1+\beta\rho)]$ (equation~\eqref{P7:eq:K_attractor}),
$\beta\rinfty=\vp$ and the golden-ratio identity $1+\vp=\vp^2$, the
$k$-essence sound speed is
\begin{equation}
  c_s^2 \;=\; 1 + \frac{\rinfty K'(\rinfty)}{K(\rinfty)}
        \;=\; 1 - \frac{\beta\rinfty}{1+\beta\rinfty}
        \;=\; 1 - \frac{\vp}{\vp^2}
        \;=\; \frac{1}{\vp^2},
        \qquad \boxed{c_s=\frac{1}{\vp}}.
  \label{P7:eq:cs_de}
\end{equation}
This is an order-unity (relativistic) speed: dark-energy perturbations
are smooth on sub-horizon scales and do not cluster, as required of a
$\Lambda$-like component. It is distinct from the $O(\beta^0)$
acoustic-metric speed $c_s^2=1$ of the acoustic-Schwarzschild section,
which sets the propagation cone of $\xi$-fluctuations on the
black-hole background: the two are the same $k$-essence sound speed
evaluated at low density and at the attractor respectively.
\end{remark}
\status{published}
\source{P7:rem:cs_de}

\begin{proposition}[EFT self-consistency and UV naturalness]
\label{P7:prop:eft_naturalness}
The effective field theory defined by the parent action~\eqref{P7:eq:parent_action}
is self-consistent for all energies $E < \LUV$: every loop correction
from matter running in the condensate background is suppressed by
$(E/\LUV)^2$ relative to the leading term, so no prediction of this
paper requires input from a UV completion above $\LUV$.
Specifically:
\begin{enumerate}[noitemsep,label=\emph{(\roman*)}]
\item \textbf{Inflationary sector.} The inflationary potential is
  evaluated at $\varphi_c\lesssim\MPl$ (Section~\ref{sec:inflation}),
  well below $\LUV\approx0.0557\,\MPl$ in energy per degree of freedom;
  corrections are $O(H_{\mathrm{inf}}^2/\LUV^2)\sim10^{-7}$.
\item \textbf{Dark energy sector.} The dark energy mass
  $m_\xi\approx H_0\sim10^{-33}\,\mathrm{eV}$ satisfies
  $m_\xi\ll\LUV$ by 52 orders of magnitude; the attractor
  solution $\rinfty=\vp/\beta$ and the thawing relation
  $w_a=-3(1+w_0)$ are independent of the UV completion.
\item \textbf{Gravitational sector.} The Planck mass is fixed by the
  density-feedback flux, $G_N=G_{\mathrm{src}}^2/(4\pi\vp^6)$
  (equation~\eqref{P7:eq:GN}); the cutoff $\LUV=\MPl/\vp^6$ sits three
  tower steps below it, so no physics above $\LUV$ is needed to
  determine $\MPl$. The Planck length $\ell_{\mathrm{Pl}}=1/\MPl$ is
  emergent, not a fundamental input.
\item \textbf{Cassini and WEP sectors.} Both the Vainshtein~\cite{Vainshtein1972} screening
  radius $r_V\approx9.3\times10^6\,\mathrm{AU}$ (Theorem~\ref{P7:thm:cassini})
  and the WEP instanton suppression $\eta_A=e^{-AS_1}$
  (Section~\ref{sec:wep}) depend on $\LUV$ and $\MPl$ only through
  the combination $m_\xi^2=\vp^{20}\Lobs\MPl^4$, which is
  fixed entirely within the EFT below $\LUV$.
\end{enumerate}
\end{proposition}
\status{published}
\source{P7:prop:eft_naturalness}

\begin{remark}[Status of the UV completion]
\label{P7:rem:uv_completion_status}
Proposition~\ref{P7:prop:eft_naturalness} establishes that the framework
is UV-safe in the sense relevant for its predictions: no measurement
described in this paper probes scales above $\LUV\approx0.0557\,\MPl$.
An explicit UV completion above $\LUV$ --- whether string embedding,
asymptotic safety, or another mechanism --- would be needed to address
questions about Planck-scale scattering or the trans-Planckian problem
for inflation, but such a completion does not affect the results here.
The characteristic condensate length scale $\ell_\mathrm{cond}=1/\LUV
\approx\vp^6\,\ell_{\mathrm{Pl}}\approx17.9\,\ell_{\mathrm{Pl}}$
is super-Planckian, consistent with the EFT being valid for
$\ell>\ell_\mathrm{cond}$.
\end{remark}
\status{published}
\source{P7:rem:uv_completion_status}

\begin{corollary}[Topological origin of the Vainshtein radius]
\label{P7:cor:vainshtein_topology}
The Vainshtein radius $r_V\approx9.3\times10^6\,\mathrm{AU}$ is
determined by the Hopf charge $Q=2$ of the condensate vacuum through
the following chain, with no free parameters:
\begin{enumerate}[noitemsep,label=\emph{(\roman*)}]
\item The $Q=2$ Hopfion vacuum, established in~\cite{Paper1}, has
  topological charge $Q_{\mathrm{condensate}}=10$ fixed by the
  WZW quantum-group order $\mathrm{ord}(q_{2I})=2(k+2)|_{k=3}=10$
  (Theorem~5.9 of~\cite{Paper1}).
\item The dark-energy scalar mass is determined by $Q_{\mathrm{condensate}}$:
  $m_\xi^2 = \vp^{2Q_{\mathrm{condensate}}}\,\Lobs\,\MPl^4
  = \vp^{20}\,\Lobs\,\MPl^4$ (equation~\eqref{P7:eq:mxi}).
\item The Brans--Dicke parameter is $\oBD=\vp^{12}/(3\beta)
  = \vp^{2(Q_{\mathrm{condensate}}-Q_{\mathrm{Hopfion}}\cdot(k+2))}/(3\beta)$,
  where $\vp^{12}=(\lam)^2$ is the square of the Bogomolny parameter.
\item Substituting both into the Vainshtein formula (Theorem~\ref{P7:thm:cassini})
  yields $r_V\propto r_S^{1/3}\,\vp^{(12-20)/3}=r_S^{1/3}\,\vp^{-8/3}$,
  so $r_V$ scales as $\vp^{-8/3}$ relative to the Schwarzschild radius,
  entirely determined by the golden ratio and the condensate topological charge.
\end{enumerate}
The Vainshtein radius is therefore a topological invariant of the $Q=2$
Hopfion condensate: it is fixed by the Hopf charge $Q_{\mathrm{Hopfion}}=2$
through the quantum-group order $Q_{\mathrm{condensate}}=10$ and the
Bogomolny parameter $\lam=\vp^6$, with no additional input.
This resolves the open problem stated in Paper~I~\cite{Paper1}.
\end{corollary}
\status{published}
\source{P7:cor:vainshtein_topology}

\begin{remark}[Frame dependence of the golden-ratio scaling]
\label{P7:rem:rV_frame}
The exponent $\vp^{-8/3}=\vp^{(12-20)/3}$ carries the $\vp^{20}$ of
the \emph{Jordan-frame} scalar mass~\eqref{P7:eq:mxi} (step~(ii)). The
physical Einstein-frame mass is $\vp$-power-free ($m_\xi\sim H_0$: the
strong nonminimal coupling $1+3\vp^7$ cancels the $\vp^{20}$
curvature). Replacing the Jordan mass by the physical one already
shifts the exponent --- to $r_V\propto\vp^4$ with $\oBD$ held fixed,
and further once $\oBD$ is itself expressed in the Einstein frame ---
so the golden-ratio scaling of $r_V$ is a Jordan-frame quantity rather
than a frame-invariant statement, and the ``topological invariant''
reading holds only in that frame. This does not affect the Cassini
resolution, which is secured independently and frame-robustly by the
chameleon channel of Remark~\ref{P7:rem:chameleon}, whose
density-dependent coupling $\oBD^{\mathrm{eff}}=\oBD(1+\bst\rho)$ is
independent of $m_\xi$.
\end{remark}
\status{published}
\source{P7:rem:rV_frame}

\begin{remark}[Lorentz violation and its recovery]
\label{P7:rem:lorentz}
Axiom~A1 introduces a preferred direction $\hat{e}_r = \nabla\rho/|\nabla\rho|$
that breaks global Lorentz invariance fundamentally.  The graviton mass
$m_g^2\propto1/\vp^6(1+\beta\rho_\infty)$ is Lorentz-breaking in the same
sense.  However, at high local density $\epsilon=\beta\rho\gg1$, the
suppression factor $S_{\mathrm{eff}}=S(\theta)/(1+\beta\rho)\to0$:
the preferred direction decouples from local physics and local Lorentz
invariance is restored to precision $O((\bst\rho)^{-1})$.
At Solar System densities $(\bst\rho_\odot)^{-1}\sim10^{-6}$, so all
PPN tests recover standard GR values at that level~\cite{Paper1}.
Conversely, in cosmic voids where $\rho\ll1/\bst$, the preferred
direction is unsuppressed and Lorentz-violating signals of order
$\sim10^{-3}$ are predicted, potentially accessible to VLBI
astrometry~\cite{Paper1}.  This density-dependent recovery is the
field-theoretic analogue of the Vainshtein mechanism and is consistent
with the Cassini resolution (Theorem~\ref{P7:thm:cassini}).
\end{remark}
\status{published}
\source{P7:rem:lorentz}

\begin{remark}[Chameleon field interpretation]
\label{P7:rem:chameleon}
The density-dependent suppression of Axiom~A2,
$S_{\mathrm{eff}}(\theta,\rho) = \sin^4\!\theta/[\vp^6(1+\beta\rho)]$,
is precisely the chameleon mechanism~\cite{KhouryWeltman2004}: the scalar $\rho$
acquires a density-dependent effective coupling
$\beta_{\mathrm{eff}}(\rho) = \beta/(1+\beta\rho)$
that runs from the UV value $\beta$ at vacuum density down to
$\beta_{\mathrm{eff}}\approx1/\rho$ in dense environments ($\beta\rho\gg1$).

This single mechanism simultaneously explains three features of the framework:
\begin{enumerate}[noitemsep,label=\emph{(\roman*)}]
  \item \textbf{Cassini resolution.}  The screened Brans-Dicke parameter
    in a dense region becomes $\oBD^{\mathrm{eff}} = \oBD(1+\beta\rho_\odot)$,
    enhanced over the vacuum value.  For $\beta\rho_\odot\gg1$ this is
    arbitrarily large, satisfying the Cassini constraint without requiring
    $\beta$ to be small; the Vainshtein~\cite{Vainshtein1972} resolution of
    Theorem~\ref{P7:thm:cassini} is complementary and operates at
    cosmological scales.
  \item \textbf{Local Lorentz recovery.}  At high density
    $S_{\mathrm{eff}}\to0$, so the preferred direction $\hat{e}_r$
    decouples from local physics; this is precisely the high-$\rho$
    recovery described in Remark~\ref{P7:rem:lorentz}.
  \item \textbf{Self-consistency.}  In~\cite{Paper1} the feedback
    functional $\Jfb = \int\mathrm{kern}/(1+\bst\,\mathrm{kern})\,\mathrm{dvol}$
    is noted to be a Born--Infeld saturation and chameleon screening
    simultaneously (Remark~4.3 of~\cite{Paper1}).
    Here the same structure operates at the cosmological level:
    the chameleon effective coupling $\beta_{\mathrm{eff}}(\rho)$
    connects the condensate-scale feedback ($\bst\approx0.452$) to
    the gravitational screening at Solar System and cosmological scales.
\end{enumerate}
The chameleon and Vainshtein mechanisms are therefore not competing
explanations of the Cassini resolution: the chameleon accounts for the
density-dependent coupling suppression in the nonlinear regime at
$r\ll r_V$, while the Vainshtein~\cite{Vainshtein1972} radius
(equation~\eqref{P7:eq:rV_exact}) sets the
scale at which nonlinear terms begin to dominate.
\end{remark}
\status{published}
\source{P7:rem:chameleon}

\begin{equation}
  ds_{\mathrm{ac}}^2 = \rho_0(r)\cdot ds_{\mathrm{Schwarzschild}}^2,
  \label{P7:eq:acoustic_schwarz}
\end{equation}

\begin{equation}
  c_s^2 = \frac{\partial P}{\partial\rho}\bigg|_{\rho_0}
        = 1 + O(\beta^3)
  \label{P7:eq:cs_schwarz}
\end{equation}

\begin{equation}
  \delta c_s^2 = -\frac{\beta^3}{4(1+\vp)} \approx -4.2\times10^{-5}
  \quad\text{at }r\sim r_s.
  \label{P7:eq:delta_cs}
\end{equation}

\begin{equation}
  \frac{dw}{da}\bigg|_{a=1}
  \;=\; \frac{d(1+w)}{da}\bigg|_{a=1}
  \;=\; 6(1+w_0).
  \label{P7:eq:dwa_slope}
\end{equation}

%════════════════════════════════════════════════════════════════════
% Dark Matter: Elastic Strain of the Condensate Director — Paper VII
%════════════════════════════════════════════════════════════════════

The framework introduces no dark-matter particle. Instead, the
condensate scalar field $\rho$ acts as the dark-matter component
through its anisotropic, nematic-like medium.  Its directional
suppression $S_{\mathrm{eff}}=\sin^4\!\theta/[\vp^6(1+\bst\rho)]$ is the
angular structure of a semi-Dirac dispersion,
$H=A\,k_\perp^2\,\sigma_x+v\,k_\parallel\,\sigma_y$, giving
\begin{equation}
  E^2 = A^2 k^4\sin^4\!\theta + v^2 k^2\cos^2\!\theta,
  \label{P7:eq:semidirac}
\end{equation}
whose heavy band carries exactly the angular factor $\sin^4\!\theta$ of
$S_{\mathrm{eff}}$. The dark-matter physics lives in the orientation
(director) field $\hat n$, not in the scalar magnitude $\rho$.
\status{published}
\source{P7:sec:dark_matter}

Visible matter imposes orientation boundary conditions on the
director; the director distorts, and the Frank elastic energy
\begin{equation}
  F=\tfrac12\int\bigl[K_{\mathrm{splay}}(\nabla\!\cdot\hat{n})^2
    +K_{\mathrm{bend}}(\hat{n}\times(\nabla\times\hat{n}))^2\bigr]\,d^3x
  \label{P7:eq:frank}
\end{equation}
of the distortion is real energy: it gravitates but carries no
electromagnetic coupling, so it is invisible. Dark matter is the
medium's elastic response to visible matter, sourced by and tracking
the baryons.
\status{published}
\source{P7:sec:director_strain}

The two director stiffnesses have opposite signs: bend is stable,
$K_{\mathrm{bend}}>0$, while splay is unstable, $K_{\mathrm{splay}}<0$.
The uniform director therefore spontaneously develops a radial splay
texture --- a disclination configuration that forms without external
tuning, the self-organised dark-matter configuration.
\status{published}
\source{P7:sec:splay_instability}

The orientation Goldstone mode has speed
\begin{equation}
  c_{\mathrm{dir}}^2 \;=\; \frac{K}{\chi},
  \label{P7:eq:cdir}
\end{equation}
where $\chi$ is the static susceptibility of the generator of director
rotations. In the semi-Dirac medium the rotational susceptibility
diverges faster with the ultraviolet scale than the Frank stiffness
$K$, so the director is heavy: $c_{\mathrm{dir}}\to0$. The director
mode is cold --- the decisive contrast with the scalar-magnitude
perturbation, whose sound speed $c_s=1/\vp$ (Remark~\ref{P7:rem:cs_de})
is relativistic and belongs to the dark-energy sector, not dark
matter. Being cold and pressureless, the director strain has no Jeans
obstruction and clusters on all cosmological and galactic scales. In
the splay channel, $K_{\mathrm{splay}}<0$ turns~\eqref{P7:eq:cdir} into
an instability, $c_{\mathrm{dir}}^2<0$, the dynamical form of the
splay self-organisation above.
\status{published}
\source{P7:sec:cold_director}

The stiffnesses have standard (gradient-squared) Frank form, so a
disclination winding gives director angle $\theta(r)\sim\ln r$, hence
$|\nabla\hat n|\sim1/r$, elastic energy density $u\sim K/r^2$, and
enclosed strain mass
\begin{equation}
  M_{\mathrm{enc}}(r)=\int u\,dV \;\propto\; r,
  \label{P7:eq:Menc}
\end{equation}
exactly $v^2=GM_{\mathrm{enc}}/r=\text{const}$ --- flat rotation
curves, with no fitted halo. A smooth, defect-free director gives the
fast-falling monopole (Keplerian); the flat curve is the signature of
the splay-driven disclination.
\status{published}
\source{P7:sec:flat_curves}

\begin{remark}[Director strain as a heavy elastic solid]
\label{P7:rem:elastic_solid}
Because the stiffnesses have standard Frank form, the director strain
is an elastic medium with genuine shear rigidity, not a collisionless
gas. It is heavy: the static rigidity is large while the shear-wave
speed stays cold ($c_{\mathrm{dir}}\to0$, equation~\eqref{P7:eq:cdir}),
the rotational inertia outgrowing the stiffness. Whether this rigidity
acts dynamically depends on the cutoff. Because $c_{\mathrm{dir}}$
falls faster than the stiffness grows, at a very large ultraviolet
cutoff the shear waves are too slow to respond on galactic timescales:
the strain is then a static, cold, pressureless gravitating field ---
phenomenologically indistinguishable from standard collisionless cold
dark matter. Only if the relevant cutoff sits near the acceleration
scale, where $c_{\mathrm{dir}}$ exceeds galactic rotation velocities,
does the elastic response act on rotation timescales; there the
medium clusters coherently and tends to suppress the small-scale
substructure a collisionless halo develops, and the splay-unstable
channel ($K_{\mathrm{splay}}<0$) renders disk warps, corrugations and
rings intrinsic elastic modes rather than tidal accidents. Robust are
the signs: shear rigidity, the coldness $c_{\mathrm{dir}}\to0$, and
the splay instability. Whether the distinguishing elastic dynamics
survive at all, and on what scales, is set by the same condensate
cutoff as the acceleration scale and is not fixed here. The overall
magnitude, by contrast, is condensate-derived: the Frank constant
follows from the condensate energy density $\rho_{\mathrm{cond}}\sim10\Lcond^4$
(fixed by $\TCMB$) and the ultralight director mass $m_\xi\sim H_0$
(fixed by the dark-energy sector), and reproduces the value required
by galactic rotation curves to within about an order of magnitude
with no galactic input.
\end{remark}
\status{published}
\source{P7:rem:elastic_solid}

The strain magnitude is set by the physical ultraviolet cutoff of the
semi-Dirac dispersion --- the condensate scale $m_\xi\sim H_0$. The
acceleration at which the strain contribution overtakes the Newtonian
term is therefore of order $cH_0$, coinciding with the empirical MOND
scale $a_0\approx cH_0/(2\pi)$ and tying the dark-matter and
dark-energy scales to the same condensate. Reproducing
$a_0\approx cH_0/(2\pi)$ requires the ratio of $m_\xi$ to the
semi-Dirac crossover scale to be of order $\vp^6$, the same
normalisation appearing in $S_{\mathrm{eff}}$ itself; this is a
framework-native scale, but the precise value is not derived here.
\status{published}
\source{P7:sec:accel_scale}

\begin{proposition}[$G_{\rm eff}(r)$ in the screened halo]
\label{P7:prop:Geff}
The BD scalar coupling to matter is characterised by
\begin{equation}
  \alpha_{\rm BD}
  \;=\; \frac{2\beta^2}{2\omega_{\rm BD}+3}
  \;\approx\; \frac{6\beta^3}{2\vp^{12}}
  \;\approx\; 8.5\times10^{-4}.
  \label{P7:eq:alpha_BD}
\end{equation}
The Vainshtein radius of a Milky-Way-mass galaxy
($M\sim6\times10^{10}\,M_\odot$) is
\begin{equation}
  r_V^{\rm MW}
  \;=\; r_V^\odot\times\!\left(\frac{M_{\rm MW}}{M_\odot}\right)^{\!1/3}
  \;\approx\; 0.045\,\mathrm{kpc}\times3915
  \;\approx\; 176\,\mathrm{kpc},
  \label{P7:eq:rV_MW}
\end{equation}
placing the entire stellar disk ($r\lesssim15\,\mathrm{kpc}\ll r_V^{\rm MW}$)
deep inside the Vainshtein sphere.
The effective Newton constant at flat-rotation-curve radii $r\lesssim30\,\mathrm{kpc}$
is
\begin{equation}
  \frac{G_{\rm eff}(r)}{G_N}
  \;=\; 1 + \alpha_{\rm BD}\!\left(\frac{r}{r_V^{\rm MW}}\right)^{\!3/2}
  \;\leq\; 1 + 6\times10^{-5}
  \qquad(r\leq30\,\mathrm{kpc}),
  \label{P7:eq:Geff_screened}
\end{equation}
indistinguishable from $G_N$ at the level of rotation-curve measurements.
\end{proposition}
\status{published}
\source{P7:prop:Geff}

Because the director strain is pressureless, carries no independent
sound speed, and is sourced by and comoving with the baryons, it adds
no acoustic feature of its own: it shifts no acoustic scale. This
no-shift property is firm; it does not, on its own, fix the
\emph{absolute} acoustic scale. The standard $r_s\approx147\,\mathrm{Mpc}$
is set by the pre-recombination matter budget, so recovering it
requires a cold-matter component present at recombination, which is
tied to the epoch at which the nematic medium orders (not settled
here): the self-organised texture orders only at $z\sim0$ and is
absent at recombination, while the recombination-era presence of the
adiabatic baryon-sourced strain is left open. The firm prediction is
therefore the \emph{absence of a shift}; recovery of the standard
scale is contingent on that ordering epoch.
\status{published}
\source{P7:sec:bao}

\begin{equation}
  \boxed{w_a \;=\; -3(1+w_0),}
  \label{P7:eq:wa_derived}
\end{equation}

\begin{theorem}[Thawing quintessence relation]
\label{P7:thm:thawing}
If $\rho$ has not fully frozen at the attractor ($\dot\rho_0\neq0$),
the CPL equation-of-state parameters satisfy
\begin{equation}
  \boxed{w_a \;=\; -3(1+w_0),}
  \label{P7:eq:thawing}
\end{equation}
where $w_a\equiv-dw/da\big|_{a=1}$ and $w_0=w(a=1)$.
This is an exact result at leading order in the slow-roll expansion
$\dot\rho^2\ll\vp^8\Lobs$ (i.e.\ $|1+w_0|\ll1$).
For the DESI~DR1 value $w_0=-0.727\pm0.067$, the prediction is
$w_a=-0.819$, within $0.80\sigma$ of the measured
$w_a=-1.05^{+0.31}_{-0.27}$~\cite{DESI2024}.
\end{theorem}
\status{published}
\source{P7:thm:thawing}
\begin{proof}
Near the attractor $\rinfty=\vp/\beta$, the kinetic prefactor
$K(\rho)=1/[\vp^6(1+\beta\rho)]$ evaluates to
\begin{equation}
  K(\rinfty)
  \;=\; \frac{1}{\vp^6(1+\beta\rinfty)}
  \;=\; \frac{1}{\vp^6(1+\vp)}
  \;=\; \frac{1}{\vp^6\cdot\vp^2}
  \;=\; \frac{1}{\vp^8},
  \label{P7:eq:K_attractor}
\end{equation}
using the golden-ratio identity $1+\vp=\vp^2$. With
$p_{\mathrm{DE}}=K\dot\rho^2/2-\Lobs$ and
$\rho_{\mathrm{DE}}=K\dot\rho^2/2+\Lobs$,
\begin{equation}
  1+w \;=\; \frac{K\dot\rho^2}{\Lobs} \;=\; \frac{\dot\rho^2}{\vp^8\Lobs} \;\ll\; 1.
  \label{P7:eq:1pw}
\end{equation}
The slow-roll homogeneous field equation $\ddot\rho+3H\dot\rho+\vp^8\Lobs=0$
gives, in the Hubble-friction dominated regime,
\begin{equation}
  3H\dot\rho \;\approx\; -m_\xi^2\delta\rho, \qquad \delta\rho\equiv\rinfty-\rho.
  \label{P7:eq:slow_roll}
\end{equation}
Differentiating~\eqref{P7:eq:1pw} and using $\ddot\rho\approx-3H\dot\rho$:
\begin{equation}
  \frac{\dot w}{H}
  \;=\; \frac{d\ln(1+w)}{d\ln a}
  \;=\; \frac{2\dot\rho\ddot\rho/(\vp^8\Lobs)}{(1+w)\,H}
  \;\approx\; -6(1+w),
  \label{P7:eq:dwda}
\end{equation}
so $(1+w)\propto a^{-6}$ backwards from today; with $1+w\to0$ as
$a\to0$ (frozen field), the growing solution is
$(1+w)(a)=(1+w_0)a^6$, giving the CPL slope
$dw/da|_{a=1}=6(1+w_0)$~\eqref{P7:eq:dwa_slope}, and, matching
$w(a)=w_0+w_a(1-a)$ over the growth-mode mean, $w_a=-3(1+w_0)$.
\end{proof}

\begin{remark}[DESI consistency and the frozen-field tension]
\label{P7:rem:desi_consistency}
The DESI~DR1 central values $w_0=-0.727\pm0.067$, $w_a=-1.05^{+0.31}_{-0.27}$
give $-3(1+w_0)=-0.819$ versus $w_a^\mathrm{DESI}=-1.05$: a $0.80\sigma$
agreement.  The frozen-field limit $w_0=-1$, $w_a=0$ is $\approx4\sigma$ from
the DESI central values.
If DESI~DR2 confirms $w_0\neq-1$ at $>3\sigma$, the framework
accommodates this via a nonzero initial displacement
$\delta\rho_0=\rinfty-\rho_\mathrm{today}\sim\sqrt{\Lobs}/\MPl^2$,
a natural scale with no fine-tuning.
The one-parameter prediction $w_a=-3(1+w_0)$ is falsifiable:
any measurement of $w_0$ uniquely pins $w_a$.
\end{remark}
\status{published}
\source{P7:rem:desi_consistency}

